\section{October 5, 2026}

Throughout, $V$ is a vector space over $\bb{F}$, and
$\mathcal L(V)$ denotes the space of linear operators on $V$.
We use $\operatorname{Id}_V$ for the identity operator and
$\mathcal P(\bb{F})$ for the space of polynomials over $\bb{F}$.

\subsection{The number of distinct eigenvalues}

Recall Corollary~\ref{cor:eigenvalue-count-bound}: an operator on a
finite-dimensional vector space $V$ has at most $\dim V$ distinct
eigenvalues. Indeed, if $\lambda_1,\ldots,\lambda_m$ are distinct
eigenvalues of $T$ and $v_1,\ldots,v_m$ are corresponding eigenvectors,
then Proposition~\ref{prop:distinct-eigenvalues-independent} says that
these vectors are linearly independent. Hence $m\leq\dim V$.

\subsection{Powers of an operator}

\begin{definition}[Powers of an operator]
  Let $T\in\mathcal L(V)$. For a positive integer $m$, define
  \[
    T^m=\underbrace{T\circ\cdots\circ T}_{m\text{ times}},
    \qquad T^0=\operatorname{Id}_V.
  \]
  If $T$ is invertible, also define
  \[
    T^{-m}=(T^{-1})^m\qquad(m\geq1).
  \]
\end{definition}

Composition of operators will often be written as multiplication:
$ST=S\circ T$, so $(ST)(v)=S(T(v))$. Associativity of composition gives
\[
  T^mT^n=T^{m+n},\qquad (T^m)^n=T^{mn}.
\]
These identities hold for nonnegative integers $m,n$ for every
operator $T$, and for all integers $m,n$ when $T$ is invertible.

\subsection{Polynomials applied to operators}

\begin{definition}[Polynomial of an operator]
  If $T\in\mathcal L(V)$ and
  \[
    p(z)=a_0+a_1z+\cdots+a_mz^m\in\mathcal P(\bb{F}),
  \]
  define the operator
  \[
    p(T)=a_0\operatorname{Id}_V+a_1T+\cdots+a_mT^m
      =\sum_{j=0}^m a_jT^j.
  \]
\end{definition}

The constant term becomes a scalar multiple of the identity operator.
Thus $p(T)$ is an operator on $V$, and $p(T)v$ is a vector in $V$.

\begin{example}[A polynomial in the differentiation operator]
  Let $D\in\mathcal L(\mathcal P(\bb{F}))$ be the differentiation
  operator $Dq=q'$, and let $p(z)=7+z+z^2$. Then
  \[
    p(D)=7\operatorname{Id}_{\mathcal P(\bb{F})}+D+D^2.
  \]
  Consequently, for every $q\in\mathcal P(\bb{F})$,
  \[
    p(D)q=7q+Dq+D(Dq)=7q+q'+q''.
  \]
\end{example}

For a fixed operator $T$, the map
\[
  \mathcal P(\bb{F})\longrightarrow\mathcal L(V),
  \qquad p\longmapsto p(T),
\]
is linear. Explicitly, for $a,b\in\bb{F}$ and polynomials $p,q$,
\[
  (ap+bq)(T)=ap(T)+bq(T).
\]
This follows by adding coefficients of equal powers, inserting zero
coefficients if the degrees differ.

The product of two polynomials is defined by
$(pq)(z)=p(z)q(z)$. Polynomial evaluation at a fixed operator also
preserves products.

\begin{proposition}\label{prop:operator-polynomial-products}
  Suppose $p,q\in\mathcal P(\bb{F})$ and $T\in\mathcal L(V)$. Then
  \begin{enumerate}[label=(\roman*)]
    \item $(pq)(T)=p(T)q(T)$;
    \item $p(T)q(T)=q(T)p(T)$.
  \end{enumerate}
\end{proposition}

\begin{proof}[Proof (left as an exercise in lecture)]
  Write
  \[
    p(z)=\sum_{j=0}^m a_jz^j,\qquad
    q(z)=\sum_{k=0}^n b_kz^k.
  \]
  Distributivity and the power law give
  \[
    p(T)q(T)=\sum_{j=0}^m\sum_{k=0}^n a_jb_kT^{j+k}
      =(pq)(T).
  \]
  Since $pq=qp$, the first identity also yields
  \[
    p(T)q(T)=(pq)(T)=(qp)(T)=q(T)p(T).\qedhere
  \]
\end{proof}

\begin{remark}
  Arbitrary operators need not commute. The proposition says that
  polynomials in the same operator commute. In particular, taking
  $q(z)=z$ shows that $p(T)T=Tp(T)$.
\end{remark}

\subsection{Invariant subspaces from polynomials}

\begin{proposition}\label{prop:polynomial-kernel-image-invariant}
  If $T\in\mathcal L(V)$ and $p\in\mathcal P(\bb{F})$, then
  $\ker p(T)$ and $\operatorname{im}p(T)$ are invariant under $T$.
\end{proposition}

\begin{proof}
  First, let $u\in\ker p(T)$, so $p(T)u=0$. Because $T$ commutes
  with $p(T)$,
  \[
    p(T)(Tu)=T(p(T)u)=T(0)=0.
  \]
  Thus $Tu\in\ker p(T)$, proving invariance of the kernel.

  Next, let $u\in\operatorname{im}p(T)$. There is a vector $v\in V$
  such that $u=p(T)v$. Again using commutativity,
  \[
    Tu=T(p(T)v)=p(T)(Tv)\in\operatorname{im}p(T).
  \]
  Thus the image is invariant as well.
\end{proof}

\subsection{Existence of eigenvalues over the complex numbers}

\begin{theorem}\label{thm:complex-operator-has-eigenvalue}
  Every operator on a nonzero finite-dimensional complex vector
  space has an eigenvalue.
\end{theorem}

\begin{proof}
  Let $T\in\mathcal L(V)$, where $V$ is a complex vector space of
  dimension $n>0$. Choose $v\in V\setminus\{0\}$. The list
  \[
    v,Tv,T^2v,\ldots,T^nv
  \]
  contains $n+1$ vectors in an $n$-dimensional space, so it is
  linearly dependent. Thus there are scalars $a_0,\ldots,a_n$, not
  all zero, such that
  \[
    a_0v+a_1Tv+\cdots+a_nT^nv=0.
  \]
  Equivalently, some nonzero polynomial $p$ satisfies $p(T)v=0$.
  Choose such a polynomial of smallest degree. Its degree is at
  least one, because a nonzero constant polynomial gives
  $p(T)v=a_0v\neq0$.

  By the fundamental theorem of algebra
  (Theorem~\ref{thm:complex-polynomial-factorization}), $p$ has a root
  $\lambda\in\bb{C}$. The factor theorem gives
  \[
    p(z)=(z-\lambda)q(z),\qquad \deg q=\deg p-1.
  \]
  By Proposition~\ref{prop:operator-polynomial-products},
  \[
    p(T)=(T-\lambda\operatorname{Id}_V)q(T).
  \]
  Minimality of the degree of $p$ implies $q(T)v\neq0$, since $q$
  is a nonzero polynomial of smaller degree. Setting $w=q(T)v$, we
  therefore have
  \[
    w\neq0,\qquad
    (T-\lambda\operatorname{Id}_V)w=p(T)v=0.
  \]
  Hence $Tw=\lambda w$, so $\lambda$ is an eigenvalue of $T$ with
  eigenvector $w$.
\end{proof}

\begin{remark}[The minimal-degree polynomial used in the proof]
  The polynomial above is chosen to satisfy $p(T)v=0$ for one fixed
  nonzero vector $v$. This does not assert $p(T)=0$ as an operator
  on all of $V$. The minimal polynomial of an operator instead
  concerns an identity on the whole space. The eigenvalue argument
  here needs only the relation for $v$.
\end{remark}

\subsection{Why the hypotheses matter}

\begin{example}[A real operator without eigenvalues]
  Recall the rotation $T(x,y)=(-y,x)$ on $\bb{R}^2$ from
  Example~\ref{ex:rotation-eigenvalues} has no real eigenvalues.
  Indeed, $T^2=-\operatorname{Id}_{\bb{R}^2}$, so $Tv=\lambda v$
  for a nonzero real vector would imply
  \[
    -v=T^2v=\lambda^2v,\qquad\lambda^2=-1,
  \]
  which is impossible for $\lambda\in\bb{R}$.
\end{example}

\begin{example}[An operator on an infinite-dimensional space]
  Consider the complex vector space $\mathcal P(\bb{C})$, which
  is infinite-dimensional. Define the multiplication operator
  \[
    T\in\mathcal L(\mathcal P(\bb{C})),\qquad (Tp)(z)=zp(z).
  \]
  If $p\neq0$, then $\deg(Tp)=\deg p+1$. For $\lambda\neq0$,
  the polynomial $\lambda p$ has degree $\deg p$, so $Tp$ cannot
  equal $\lambda p$. For $\lambda=0$, the equation is also
  impossible because $Tp\neq0$. Thus $T$ has no eigenvalues.
\end{example}
